\documentclass[12pt]{article}
\usepackage{amsmath,amssymb,amsfonts,amsthm}
\usepackage[margin=1in]{geometry}

\begin{document}

\begin{center}
{\Large\bfseries Chapter 9, Problem 20}\\[6pt]
Byron \& Fuller, \textit{Mathematics of Classical and Quantum Physics}
\end{center}

\bigskip

\textbf{Problem.} Extend the Fredholm method to square-integrable kernels by noting that:

(a) We can always modify $k(x,s)$ so that $\tilde{k}(x,a) = 0$ on the diagonal without changing solutions.

(b) Define a modified Fredholm determinant $\tilde{D}(\lambda)$ by replacing diagonal entries $k(x_i,x_i)$ with zero.

\bigskip
\textbf{Solution.}

\subsection*{(a) Modifying the kernel on the diagonal}

The difficulty with extending the Fredholm method to square-integrable (but not continuous) kernels is that terms like $\int k(x,x)\,dx$ may not exist---the diagonal $\{(x,x)\}$ is a set of measure zero in $[a,b]^2$, and modifying a function on a set of measure zero does not change its $L_2$ properties.

By Problem~9.2, modifying $k(x,t)$ on a set of measure zero does not change the integral operator $K$. The diagonal $\{(x,t) : x = t\}$ is a one-dimensional curve in $[a,b]^2$ and has two-dimensional Lebesgue measure zero. Therefore, we can define:
\[
\tilde{k}(x,t) = \begin{cases} k(x,t) & \text{if } x \neq t, \\ 0 & \text{if } x = t, \end{cases}
\]
and the integral operator with kernel $\tilde{k}$ is identical to $K$ on $L_2[a,b]$.

All solutions of $f = g + \lambda Kf$ are unchanged (as $L_2$ functions), so there is no loss in assuming $k(x,x) = 0$.

\subsection*{(b) Modified Fredholm determinant}

With $\tilde{k}(x,x) = 0$, define the modified Fredholm determinant by:
\[
\tilde{D}(\lambda) = \sum_{n=0}^{\infty} \frac{(-\lambda)^n}{n!}\,\tilde{d}_n,
\]
where $\tilde{d}_n$ is the same integral as $d_n$ but with the determinant $\tilde{\Delta}_n$ obtained by replacing each diagonal entry $k(x_i, x_i)$ by zero:
\[
\tilde{\Delta}_n = \begin{vmatrix} 0 & k(x_1,x_2) & \cdots & k(x_1,x_n) \\ k(x_2,x_1) & 0 & \cdots & k(x_2,x_n) \\ \vdots & \vdots & \ddots & \vdots \\ k(x_n,x_1) & k(x_n,x_2) & \cdots & 0 \end{vmatrix}.
\]

The crucial point is that $\tilde{\Delta}_n$ involves only \emph{off-diagonal} values of $k$, so the integrals $\tilde{d}_n = \int\cdots\int \tilde{\Delta}_n\,dx_1\cdots dx_n$ are well-defined for any $L_2$ kernel. Hadamard's inequality still applies:
\[
|\tilde{\Delta}_n|^2 \le \prod_{i=1}^n \sum_{j \neq i} |k(x_i, x_j)|^2,
\]
and the convergence proof of Problem~9.13 carries through with the $L_2$ norm $\|k\|_{L_2}^2 = \int\!\!\int |k(x,t)|^2\,dx\,dt$ replacing the supremum norm.

Similarly, the modified numerator $\tilde{N}(x,y;\lambda)$ is defined with zeros on the diagonal in the appropriate places. The resolvent $\tilde{R} = \tilde{N}/\tilde{D}$ satisfies the same equations as before (Problems~9.14--9.16), and provides the solution to $f = g + \lambda Kf$ for all $\lambda$ where $\tilde{D}(\lambda) \neq 0$.

This extension, due essentially to Carleman, allows the Fredholm theory to handle all square-integrable kernels, including Green's function kernels of the form $G(x,t)/w(t)$ which may have singularities on the diagonal. \hfill $\blacksquare$

\end{document}
